%% begin commment
\iffalse
基本的な例です．
本文で用いるコマンド，環境はjbook.clsに準拠してください．
A5判の場合

\documentclass[a5/13Q]{kbdbook6}

B5判の場合
\documentclass[b5/13Q]{kbdbook6}
uplatexを使う場合は，それぞれのoptionに[uplatex]を追記します．
\documentclass[a5/13Q,uplatex]{kbdbook6}

\usepackage[dvipdfmx]{graphicx}  %% 図版処理にはdvipdfmxを使用します．
                                %% 現在では最終成果物DVIファイルではなくPDF形式が標準となっております．
\usepackage[TS1,OT1,T1]{fontenc} %% 使用するエンコーディングはT1をデフォルト
\usepackage{textcomp}%% LaTeXで使える記号類を拡張
\usepackage{amsmath,amssymb}%% 数式でamsmath使う場合．
\usepackage{makeidx}
\makeindex

各文書は章ごとにファイルを分割して\includeで読み込む形で執筆されることをお薦めいたします．
\includeonlyコマンドは各ファイルの読み込みを制御して，コンパイル時間を短縮できます．
全文書などコンパイルが必要が無い場合は，該当する章をコメントアウトします．
その場合でも，目次や章番号，章を跨いだ相互参照などの情報は確保します．

\includeonly{%%
	preface,% まえがき
	chap01,	%1章
	chap02,	%2章
	chap03,	%3章
	chap04, %4章
	chap05,	%5章
}
chap02のみ処理したい場合
\includeonly{%%
%	preface,% まえがき
	chap01,	%1章
%	chap02,	%2章
%	chap03,	%3章
%	chap04, %4章
%	chap05,	%5章
}

\includeonly{%%
	preface,% まえがき
	chap01,	%1章
	chap02,	%2章
	chap03,	%3章
	chap04, %4章
	chap05,	%5章
}
\fi
%% end commment


%LaTeX sample start
\documentclass[b5/13Q]{kbdbook6}
\usepackage[dvipdfmx]{graphicx}
\usepackage[TS1,OT1,T1]{fontenc}
\usepackage{textcomp}
\usepackage{amsmath,amssymb}
\usepackage{makeidx}
\makeindex

\includeonly{%%
	preface,% まえがき
	chap01,	%1章
%	chap02,	%2章
%	chap03,	%3章
%	chap04, %4章
%	chap05,	%5章
}

\begin{document}
\frontmatter%% 前付の指定
\chapter*{まえがき}
\markboth{まえがき}{まえがき}
□□□□□□□□□□□□□□□□□□□□□
□□□□□□□□□□□□□□□□□□□□□
□□□□□□□□□□□□□□□□□□□□□
□□□□□□□□□□□□□□□□□□□□□
□□□□□□□□□□□□□□□□□□□□□
□□□□□□□□□□□□□□□□□□□□□
□□□□□□□□□□□□□□□□□□□□□
□□□□□□□□□□□□□□□□□□□□□
□□□□□□□□□□□□□□□□□□□□□
□□□□□□□□□□□□□□□□□□□□□
□□□□□□□□□□□□□□□□□□□□□
□□□□□□□□□□□□□□□□□□□□□
□□□□□□□□□□□□□□□□□□□□□
□□□□□□□□□□□□□□□□□□□□□
□□□□□□□□□□□□□□□□□□□□□
□□□□□□□□□□□□□□□□□□□□□
□□□□□□□□□□□□□□□□□□□□□
□□□□□□□□□□□□□□□□□□□□□
□□□□□□□□□□□□□□□□□□□□□
□□□□□□□□□□□□□□□□□□□□□
□□□□□□□□□□□□□□□□□□□□□
□□□□□□□□□□□□□□□□□□□□□
□□□□□□□□□□□□□□□□□□□□□
□□□□□□□□□□□□□□□□□□□□□
□□□□□□□□□□□□□□□□□□□□□
□□□□□□□□□□□□□□□□□□□□□
□□□□□□□□□□□□□□□□□□□□□
□□□□□□□□□□□□□□□□□□□□□
□□□□□□□□□□□□□□□□□□□□□
□□□□□□□□□□□□□□□□□□□□□
□□□□□□□□□□□□□□□□□□□□□
□□□□□□□□□□□□□□□□□□□□□
□□□□□□□□□□□□□□□□□□□□□
□□□□□□□□□□□□□□□□□□□□□
□□□□□□□□□□□□□□□□□□□□□
□□□□□□□□□□□□□□□□□□□□□
□□□□□□□□□□□□□□□□□□□□□
□□□□□□□□□□□□□□□□□□□□□


\vspace{\baselineskip}
2014年○月

\hspace*{\fill}著者名

%% まえがきを記述したTeXファイル
\tableofcontents
\mainmatter%% 本文開始指定
%% 本文の記述はjbook準拠としてください．
\include{chap01}%% 1章を記述したTeXファイル（以下同）
\include{chap02}
\include{chap03}
\include{chap04}
\include{chap05}
\backmatter
\printindex
\end{document}
